\section{Classical structure and the additional quantifiers}
Classical stochastic-semigroup and nonnegative-matrix-group structure,
including Flor \cite{Flor} and Schwarz \cite{Schwarz58}, supplies recurrent
idempotents and permutation actions on their finite recurrent simplexes.
The purification proof uses that structure inside the joint physical and
stochastic closure. The extra conclusion keeps the prescribed closed
numerical error balls and the original label budget, then descends the
hidden action to the physical group.

Fijalkow--Paperman \cite[Theorems~1 and 10]{FP} characterize advice
regular languages and prove neutral-letter collapse. In fact their
Myhill--Nerode proof preserves the relevant deterministic bound $K$;
we do not claim that preservation of a state bound, by itself, is new.
The difference here is the stochastic simplex, common closed numerical
error, approximate executable returns and an occupation obstruction with
unrestricted intermediate widths. Our proof keeps the homogeneous
neutral table inside the terminal decoder rather than identifying it
with a transition identity.

Rational and weighted-series minimization \cite{BR,Paz} concerns linear
representations and Hankel spaces. Positive realization \cite{BF04}
adds a cone constraint. Neither linear dimension nor a separate
nonnegative factorization at each cut is used here as a substitute for
one compatible stochastic realization. Shitov's exact complexity
statement \cite[Theorem~2]{Shitov} is imported in the finite-group
hardness proof, not re-proved as a realization theorem.

For group and quantum automata, Brodsky--Pippenger
\cite[Theorems~3.3 and 3.6]{BP} compare cut-point languages, whereas our
physical action retains a full continuous numerical interface and its
closed error tolerance. Distribution metrics \cite{FSZ} and topological
recognition \cite{Steinberg} give adjacent equivalence and dynamical
frameworks; they are not assumed to identify the same-width stochastic
minimum. Classical hidden Markov realization \cite{Heller58} imposes
process consistency that is not requested by a single terminal query.
No online or general irreversible-group-quotient conclusion follows.

The effectivity assertions must be kept separate. A finite continuous
quotient may exist without an input encoding. Algebraic closure
algorithms permit computation for rational matrix input. An explicitly
tabulated finite group has the displayed existential-real decision
classification. Constructing a minimum realization and compiling a
supplied encoded realization are further tasks with different costs.
The finite-precision output sizes are not, in general, runtime estimates.
This theorem-level comparison is author-side positioning, not independent
priority clearance.

\subsection{Controlled stochastic processes}
Singh, James and Rudary \cite[Theorems~2 and~6]{SJR2004} distinguish
finite hidden-state dimension from linear predictive dimension. Their
first theorem bounds the rank of the system-dynamics matrix by the
number of POMDP states; their second characterizes finite linear
predictive representations, whose update coefficients need not be
nonnegative. Neither rank counts distinct future rows.
Our causal theorem instead uses repeated nondisturbing probes to make
finitely many selected future responses statistically distinguishable
by one common experiment. It then tests every horizon-dependent
nonnegative instrument across the same register. The new condition is
uniform total variation of whole adaptive transcripts. It is not
inferred from a linear-rank calculation and does not claim a general
HMM classification. The proof supplies the stated result, not an
independent certification of priority across the realization literature.
